\chapter{Cơ sở Lý thuyết và Tổng quan Nghiên cứu}
\label{ch:theoretical_background}

Chương này trình bày các nguyên lý lý thuyết nền tảng của luận văn, bao gồm: thực trạng vận hành SOC và nút thắt nhận thức, cơ sở luật tất định, thống kê trực tuyến và học máy Tầng 1, kiến trúc tác tử và mô hình ngôn ngữ lớn Tầng 2, nguy cơ đối kháng và kiểm toán mật mã, cùng tổng quan các nghiên cứu liên quan nhằm xác định khoảng trống nghiên cứu mà SENTINEL giải quyết.

\section{Thực trạng Vận hành SOC và Nút thắt Nhận thức}

Mục này đi sâu phân tích các nguyên nhân kiến trúc cốt lõi gây nên tình trạng quá tải cảnh báo đã được phác thảo tại Chương~\ref{ch:introduction}. Trong mô hình SOC tiêu chuẩn, công tác giám sát dựa vào việc hợp nhất nhật ký thô từ các cảm biến như Tường lửa thế hệ mới (NGFW), Hệ thống phát hiện xâm nhập (IDS/IPS) và giải pháp Phản hồi Điểm cuối (EDR) về nền tảng SIEM tập trung~\cite{siem2021}, với khối lượng vượt xa năng lực phân tích thủ công của con người~\cite{bigdata2019}. Nút thắt khởi nguồn từ cơ chế vận hành của các cảm biến đầu vào: chúng thực hiện phán quyết chủ yếu dựa trên chữ ký tĩnh hoặc các luật tương quan tất định. Do áp lực tránh bỏ lọt sự cố luôn đi kèm thiệt hại lớn, người vận hành bị buộc phải hạ ngưỡng phát hiện. Tuy nhiên, việc hạ ngưỡng trong các hệ thống dựa trên chữ ký chỉ làm gia tăng khối lượng báo động giả gửi về hàng đợi mà không thể nâng cao năng lực phân tích ngữ cảnh~\cite{alertfatigue2022, alahmadi2022}.

Ở lớp trên, các nền tảng Điều phối và Tự động hóa An ninh (SOAR) được triển khai nhằm xả tải thao tác bằng cách thực thi các kịch bản ứng phó (playbooks) lập trình sẵn~\cite{soar2025}. Dù vậy, ranh giới của SOAR bị siết chặt bởi tính cố định của kịch bản: một playbook chỉ kích hoạt thành công khi cảnh báo khớp chính xác với tình huống đã được lường trước. Đối với các chuỗi tấn công phức tạp, hành vi bất thường mới hoặc bằng chứng phi cấu trúc nằm rải rác, cả ba lớp SIEM, IDS/IPS và SOAR đều không thể hình thành phán quyết ngữ cảnh. Toàn bộ các ca nghi ngờ này buộc phải chuyển về hàng đợi thủ công, biến chuyên viên trực ca thành nút thắt cổ chai nhận thức duy nhất của toàn bộ hệ thống.

\section{Lý thuyết Luật tất định, Thống kê Trực tuyến và Học máy Tầng 1}
\label{sec:welford}

Yêu cầu vận hành ở tốc độ đường truyền đặt ra ràng buộc chặt chẽ về mặt thuật toán cho Tầng 1: các cơ chế xử lý phải đạt độ phức tạp thời gian và bộ nhớ hằng số $\mathcal{O}(1)$. Đầu tiên, lớp phòng thủ dựa trên luật (Rule-based) sử dụng các tập hợp đối sánh tĩnh và tra cứu cổng dịch vụ để loại bỏ tức thời các hành vi bất thường thô mà không tiêu tốn tài nguyên tính toán. Tiếp theo, cơ chế theo dõi đường nền áp dụng công thức truy hồi tăng dần Welford~\cite{welford1962} để cập nhật liên tục giá trị trung bình $\mu_k$ và tổng bình phương độ lệch $S_k$ theo từng quan sát, tính toán độ lệch chuẩn $\sigma_k = \sqrt{S_k / (k-1)}$ và điểm $Z$-score động mà không cần lưu trữ nhật ký quá khứ:
\begin{equation}
\label{eq:welford}
    \mu_k = \mu_{k-1} + \frac{x_k - \mu_{k-1}}{k}, \quad S_k = S_{k-1} + (x_k - \mu_{k-1})(x_k - \mu_k)
\end{equation}

Cuối cùng, các mô hình Học máy dựa trên Cây quyết định tăng cường độ dốc (GBDT), tiêu biểu là thuật toán LightGBM~\cite{ke2017lightgbm}, cung cấp năng lực phân loại nhanh các mẫu dữ liệu đa chiều. Việc kết hợp lọc luật tất định, chuẩn hóa thống kê $Z$-score và phân loại học máy tạo thành tuyến phòng thủ đa tầng tốc độ cao, vừa đảm bảo xả tải lưu lượng lớn vừa duy trì tính ổn định trước các biến thể dữ liệu đối kháng.

\section{Mô hình Ngôn ngữ Lớn Cục bộ và Trí tuệ Tác tử Tầng 2}
\label{sec:local_llm}

Vận hành mô hình suy luận hoàn toàn cục bộ là điều kiện bắt buộc để bảo đảm chủ quyền dữ liệu an ninh nội bộ, nhưng đồng thời đặt ra giới hạn nghiêm ngặt về tài nguyên bộ nhớ GPU. Kỹ thuật lượng tử hóa giúp nén biểu diễn trọng số sang lưới số nguyên thô hơn, đánh đổi một lượng nhỏ độ trung thực số học để giảm đáng kể dung lượng bộ nhớ~\cite{quantization2021}. Ở định dạng GGUF Q4\_K\_M, mô hình chuyên biệt an ninh Foundation-Sec-8B-Instruct~\cite{foundation_sec_instruct, foundation_sec_2024} chỉ chiếm khoảng 7–8\,GB VRAM, chừa đủ không gian cho bộ đệm ngữ cảnh và chỉ mục truy xuất trên card đồ họa GPU cục bộ.

Tầng suy luận điều phối quy trình phân tích thông qua đồ thị trạng thái LangGraph~\cite{langgraph2024}. Nhằm khắc phục hạn chế ảo giác và neo phán quyết vào chứng cứ thực tế~\cite{rag2020}, hệ thống kết hợp kỹ thuật truy xuất thông tin lai giữa tìm kiếm véc-tơ dày (FAISS)~\cite{faiss2019} và đối sánh từ khóa thưa (BM25)~\cite{bm25_2009}. Kết quả truy xuất từ hai phương pháp được hợp nhất theo thứ hạng bằng thuật toán Reciprocal Rank Fusion (RRF)~\cite{rrf2009}:
\begin{equation}
\label{eq:rrf}
    RRF\_Score(d) = \sum_{m \in M} \frac{1}{k + r_m(d)}
\end{equation}
giúp đối chiếu minh bạch dữ liệu nhật ký với các mã kỹ thuật MITRE ATT\&CK~\cite{strom2018mitre} và hướng dẫn ứng phó sự cố NIST SP 800-61r2~\cite{nist80061}.

\section{An ninh Đối kháng AI và Mật mã học Kiểm toán Nhật ký}
\label{sec:adversarial_theory}

Tích hợp Mô hình Ngôn ngữ Lớn vào quy trình phân tích an ninh làm phát sinh rủi ro bị tấn công tiêm nhiễm câu lệnh qua nhật ký (Log Prompt Injection)~\cite{promptinjection2023, owasp_llm2025}. Do nội dung nhật ký do chính bên tấn công khởi tạo, các giải pháp phòng thủ dựa trên phân loại nội dung hoặc lọc từ khóa xác suất thường dễ bị qua mặt trước các biến thể câu lệnh đối kháng mới~\cite{watchtower2026}. Căn cứ lý thuyết đó đặt ra yêu cầu phải bảo vệ tầng suy luận bằng cơ chế cách ly cấu trúc tất định (Delimited Data Encapsulation), phân tách ranh giới tuyệt đối giữa chỉ thị điều khiển và dữ liệu thô bị nghi nhiễm.

Đối với tính toàn vẹn của hồ sơ phân tích, các vết chứng cứ pháp y được niêm phong theo chuỗi bằng thuật toán mã hóa HMAC-SHA256~\cite{hmac1997} sử dụng khóa bí mật cục bộ. Về mặt ranh giới lý thuyết, cơ chế mã hóa đối xứng HMAC đảm bảo khả năng phát hiện mọi hành vi chỉnh sửa hoặc giả mạo nhật ký, nhưng không thay thế cho thuộc tính chống chối bỏ (non-repudiation) vốn đòi hỏi hạ tầng chữ ký số bất đối xứng.

\section{Tổng quan Nghiên cứu Liên quan và Khoảng trống Nghiên cứu}

Nghiên cứu ứng dụng trí tuệ nhân tạo trong vận hành an ninh mạng đã tiến triển qua ba nhóm giải pháp kiến trúc chính, theo khảo sát hệ thống của Zhang và cộng sự~\cite{llmsec_survey2025}.

Nhóm thứ nhất, phòng thủ doanh nghiệp truyền thống, gồm các nền tảng SIEM như IBM QRadar và giải pháp SOAR như Palo Alto Cortex XSOAR~\cite{siem2021, soar2025}, tương quan sự kiện bằng tập luật tĩnh và chạy kịch bản định sẵn. Bổ trợ cho nhóm này là các kỹ thuật phân tách nhật ký trực tuyến kinh điển như Drain~\cite{he2017drain}, giúp rút gọn cú pháp nhưng chưa thể phân tích ngữ cảnh hành vi. Qua phỏng vấn người vận hành SOC, AlAhmadi và cộng sự~\cite{alahmadi2022} ghi nhận chất lượng cảnh báo thấp là nguyên nhân hàng đầu gây nên tình trạng quá tải. Ưu thế cốt lõi của nhóm này là độ trễ rất thấp, đặc tính mà Tầng 1 của SENTINEL kế thừa và phát triển.

Nhóm thứ hai, tác tử phân loại có tăng cường truy xuất, tiêu biểu là CyberRAG~\cite{cyberrag2025}, điều phối các bộ phân loại tinh chỉnh theo từng họ tấn công phía sau một tác tử LLM và công bố độ chính xác trên 94\% trên các phép thử tấn công web. Công trình này giải quyết chất lượng phân tích cho những sự kiện đã chọn, nhưng thiếu tầng lọc tốc độ đường truyền ở đầu nguồn để quản lý lưu lượng lớn.

Nhóm thứ ba, đa tác tử và chú trọng quản trị, gồm AutoBnB-RAG~\cite{autobnb2025}, LanG~\cite{lang2026} và Policy-Guided Threat Hunting~\cite{splunkllm2026}. AutoBnB-RAG đánh giá phản ứng sự cố trong môi trường mô phỏng trò chơi bàn \textit{Backdoors \& Breaches}. LanG dựng trên LangGraph với các điểm chốt duyệt thủ công và phòng thủ tiêm nhiễm hai lớp (biểu thức chính quy kết hợp Llama Prompt Guard 2 đạt F1 98,1\%). Policy-Guided Threat Hunting sử dụng bộ tự mã hoá tái dựng và học tăng cường sâu hai lớp để tiền phân loại trước khi đưa vào LLM.

Khoảng trống nghiên cứu được xác định rõ ràng: phân tầng chi phí đã xuất hiện ở Policy-Guided Threat Hunting, phòng thủ tiêm nhiễm đã có ở LanG, và vận hành cục bộ đã được triển khai ở một số công trình. Tuy nhiên, chưa có kiến trúc đơn lẻ nào hợp nhất thành công: (1) giải quyết phần lớn lưu lượng ở tầng thống kê/học máy với độ trễ dưới 1\,ms và tỷ lệ xả tải luôn đi kèm tỷ lệ tấn công nền của luồng dùng để đo; (2) suy luận có trạng thái chạy trên GPU cục bộ dưới quy tắc neo truy xuất nghiêm ngặt; và (3) bảo chứng phòng thủ tiêm nhiễm dựa trên cấu trúc kết hợp với niêm phong pháp y HMAC-SHA256. Bảng~\ref{tab:related_work_comparison} đối chiếu chi tiết SENTINEL với các nhóm giải pháp đại diện.

\begin{table}[htbp]
\centering
\caption[So sánh đối chiếu với các giải pháp hiện hành]{Bảng so sánh đối chiếu giữa SENTINEL và các giải pháp phòng thủ hiện hành}
\label{tab:related_work_comparison}
\footnotesize
\setlength{\tabcolsep}{4pt}
\begin{tabularx}{\textwidth}{>{\raggedright\arraybackslash}p{2.6cm} >{\raggedright\arraybackslash}X >{\raggedright\arraybackslash}X >{\raggedright\arraybackslash}X}
\toprule
\textbf{Tiêu chí Đánh giá} & \textbf{SIEM / SOAR (QRadar, XSOAR)} & \textbf{Hệ tác tử LLM (CyberRAG, LanG, Policy-Guided)} & \textbf{Kiến trúc SENTINEL (Luận văn)} \\
\midrule
Kỹ thuật Cốt lõi & Tập luật tương quan tĩnh, kịch bản Python/JS cứng & RAG tác tử, điều phối đa tác tử, phân loại sơ bộ bằng DRL & Bộ luật chữ ký Rule-based + Welford $\mathcal{O}(1)$ + LightGBM T1, LangGraph Agent T2 \\
Mục tiêu Vận hành & Gom nhật ký tập trung và phản hồi theo kịch bản & Tự động hóa điều tra sự cố bằng suy luận ngôn ngữ & Giải quyết lưu lượng tại tầng rẻ, dành tầng đắt cho ca cần suy luận \\
Độ trễ \& Phân tầng & Một tầng rẻ, không suy luận ngữ cảnh & Có ở một số công trình (DRL sơ bộ); chưa công bố tỷ lệ xả tải kèm tỷ lệ tấn công nền & Hai tầng: trung bình 0,182\,ms ở T1; xả tải 97,5\% lưu lượng trước T2 ở tỷ lệ tấn công nền 9,8\% \\
Chủ quyền Dữ liệu & Cục bộ nội bộ & Triển khai cục bộ nhưng không ấn định giới hạn phần cứng nghiêm ngặt & Air-gapped 100\% trên GPU cục bộ \\
Chống Prompt Inj. & Không áp dụng & Bộ phân loại xác suất (LanG: regex + Prompt Guard 2, F1 98,1\%) & Bảo chứng cấu trúc Đóng gói Dữ liệu Phân định, độc lập với câu chữ payload \\
Kiểm toán Pháp y & Nhật ký văn bản tiêu chuẩn (dễ bị chỉnh sửa) & Có công bố nhật ký kiểm toán (LanG: \texttt{mcp\_audit.db}); không niêm phong mật mã & Chuỗi HMAC-SHA256 phát hiện giả mạo vết chứng cứ \\
\bottomrule
\end{tabularx}
\end{table}
